\chapter{Equations in Two Variables}

\section*{12.1. Quasiconformal Mappings}

\begin{definition}[Quasiconformal mappings]
\label{gt-ch12-def-quasiconformal-mapping}
\dcref{ch12:12.1--12.2}
\group{chapter-12}\source{gilbarg-trudinger:page-0306}
Let $\Omega\subset\mathbb R^2$ and let
$w=p+iq\in C^1(\Omega;\mathbb C)$.  For $K\ge1$, the mapping is
$K$-quasiconformal when
\begin{equation}\tag{12.1}
p_x^2+p_y^2+q_x^2+q_y^2
\le2K(p_xq_y-p_yq_x)
\end{equation}
throughout $\Omega$.  More generally, it is $(K,K')$-quasiconformal, where
$K\ge1$ and $K'\ge0$, when
\begin{equation}\tag{12.2}
p_x^2+p_y^2+q_x^2+q_y^2
\le2K(p_xq_y-p_yq_x)+K'.
\end{equation}
The same definitions apply to continuous $p,q\in W^{1,2}_{\mathrm{loc}}$ by
requiring these inequalities almost everywhere.
\end{definition}

For a disk $B_r(z)\Subset\Omega$, set
\begin{equation}\tag{12.3}
\mathcal D(r;z)=\iint_{B_r(z)}|Dw|^2\,dx\,dy.
\end{equation}

\begin{lemma}[Growth of the Dirichlet integral]
\label{gt-ch12-lem-dirichlet-integral-growth}
\dcref{ch12:12.1}\group{chapter-12}\source{gilbarg-trudinger:page-0307}
Let $w=p+iq$ be $(K,K')$-quasiconformal in
$B_R=B_R(z_0)$, with $K>1$, $K'\ge0$, and $|p|\le M$.  For
$r\le R/2$,
\begin{equation}\tag{12.4}
\mathcal D(r)\le C\left(\frac rR\right)^{2\alpha},
\qquad
\alpha=K-\sqrt{K^2-1},
\qquad
C=C_1(K)(M^2+K'R^2).
\end{equation}
If $K'=0$, this remains valid for $K=1$.
\end{lemma}
\begin{proof}
\sketch Green's formula and (12.2) give
\begin{equation}\tag{12.5}
\mathcal D(r)\le
2K\int_{\partial B_r}p\,q_s\,ds+\pi K'r^2.
\end{equation}
Since $\mathcal D'(r)=\int_{\partial B_r}|Dw|^2\,ds$, Cauchy--Schwarz gives
\begin{equation}\tag{12.6}
\int_{\partial B_r}p\,q_s\,ds
\le M\sqrt{2\pi r\,\mathcal D'(r)}.
\end{equation}
Consequently,
\begin{equation}\tag{12.7}
[\mathcal D(r)-k_1]^2\le k_2r^2\mathcal D'(r),
\qquad
k_1=\pi R^2K',
\quad k_2=8\pi M^2K^2,
\end{equation}
and integration on $(R/2,R)$ yields
\begin{equation}\tag{12.8}
\mathcal D(R/2)\le\frac{8\pi}{\log2}M^2K^2+\pi R^2K'.
\end{equation}

Put $\alpha=K-\sqrt{K^2-1}$, so
$K=(1+\alpha^2)/(2\alpha)$.  The quasiconformal inequality implies, for
every unit direction $s$,
\begin{equation}\tag{12.9}
|w_s|^2\le\frac1{1+\alpha^2}
\left(|Dw|^2+\frac{2\alpha^2K'}{1-\alpha^2}\right).
\end{equation}
If $\bar p$ is the mean of $p$ on $\partial B_r$, then
\begin{equation}\tag{12.10}
\int_{\partial B_r}p\,q_s\,ds
\le\frac12\int_{\partial B_r}
\left(\frac{(p-\bar p)^2}{r}+rq_s^2\right)\,ds,
\end{equation}
while Wirtinger's inequality gives
\begin{equation}\tag{12.11}
\int_{\partial B_r}(p-\bar p)^2\,ds
\le r^2\int_{\partial B_r}p_s^2\,ds.
\end{equation}
Combining (12.5), (12.9)--(12.11) produces
\begin{equation}\tag{12.12}
\mathcal D(r)\le\frac r{2\alpha}\mathcal D'(r)+kr^2,
\qquad
k=\pi K'\left(1+\frac{2\alpha}{1-\alpha^2}\right).
\end{equation}
Integration from $r$ to $r_0$ gives
\begin{equation}\tag{12.13}
\mathcal D(r)\le
\left[\mathcal D(r_0)+\frac{\alpha}{1-\alpha}kr_0^2\right]
\left(\frac r{r_0}\right)^{2\alpha}.
\end{equation}
Use $r_0=R/2$ and (12.8).  When $K'=0$, the limiting case $K=1$ is
unchanged.
\end{proof}

\begin{lemma}[Dirichlet growth implies H\"older continuity]
\label{gt-ch12-lem-dirichlet-growth-holder}
\dcref{ch12:12.2}\group{chapter-12}\source{gilbarg-trudinger:page-0310}
Let $w\in C^1(\Omega)$ and
$\widetilde\Omega\Subset\Omega$ with
$\operatorname{dist}(\widetilde\Omega,\partial\Omega)>R$.  Suppose
\[
\mathcal D(r;z)\le Cr^{2\alpha}
\]
for every $z\in\widetilde\Omega$ and $r\le R'\le R$.  If
$z_1,z_2\in\widetilde\Omega$ and $|z_1-z_2|\le R'$, then
\[
|w(z_2)-w(z_1)|
\le2\sqrt{\frac C\alpha}\,|z_2-z_1|^\alpha.
\]
\end{lemma}
\begin{proof}
\sketch Apply the two-dimensional Morrey estimate to $w$ and use
Cauchy--Schwarz to express its hypothesis through $\mathcal D(r;z)$.
\end{proof}

\begin{theorem}[Interior H\"older estimate for quasiconformal mappings]
\label{gt-ch12-thm-holder-quasiconformal-mappings}
\dcref{ch12:12.3}\group{chapter-12}\source{gilbarg-trudinger:page-0311}
Let $w$ be $(K,K')$-quasiconformal in $\Omega$, with $K>1$, $K'\ge0$, and
$|w|\le M$.  If
$\widetilde\Omega\Subset\Omega$ and
$\operatorname{dist}(\widetilde\Omega,\partial\Omega)>d$, then
\begin{equation}\tag{12.14}
|w(z_2)-w(z_1)|
\le C\left|\frac{z_2-z_1}{d}\right|^\alpha,
\qquad
\alpha=K-\sqrt{K^2-1},
\end{equation}
for all $z_1,z_2\in\widetilde\Omega$, where
$C=C_1(K)(M+d\sqrt{K'})$.  If $K'=0$, the conclusion also holds for
$K=1$, with $C=C_1(K)M$.
\end{theorem}
\begin{proof}
\sketch If $|z_2-z_1|\le d/2$, apply
Lemmas~\ref{gt-ch12-lem-dirichlet-integral-growth} and
\ref{gt-ch12-lem-dirichlet-growth-holder} with $R=d$.  For larger
separations, use $|w(z_2)-w(z_1)|\le2M$ and absorb this into the right-hand
side of (12.14).
\end{proof}

The exponent in (12.14) is optimal for $K>1$.  For $K=1$ and $K'>0$, the
formula gives $\alpha=1$, which is false in general; nevertheless (12.14)
holds with every $\alpha<1$ because (12.2) remains valid after increasing
$K$.  If $\widetilde\Omega$ is covered by $N$ disks of diameter $d/2$, the
hypothesis $|w|\le M$ may be weakened to $|p|\le M$, with $C$ also depending
on $N$.

If $\Omega$ is $C^1$, $w\in C^1(\overline\Omega)$, and $p=0$ on
$\partial\Omega$, boundary flattening and the odd-even reflection
\[
p(\xi,-\eta)=-p(\xi,\eta),
\qquad q(\xi,-\eta)=q(\xi,\eta)
\]
give a global H\"older estimate for $w$.  If instead
$p=\widetilde p$ on $\partial\Omega$, with
$\widetilde p\in C^1(\overline\Omega)$, apply the same argument to
$p-\widetilde p$.

\section*{12.2. H\"older Gradient Estimates for Linear Equations}

Consider
\begin{equation}\tag{12.15}
Lu=au_{xx}+2bu_{xy}+cu_{yy}=f
\end{equation}
in $\Omega\subset\mathbb R^2$, with eigenvalues $\lambda,\Lambda$ satisfying
\begin{equation}\tag{12.16}
\lambda(\xi^2+\eta^2)
\le a\xi^2+2b\xi\eta+c\eta^2
\le\Lambda(\xi^2+\eta^2)
\end{equation}
and
\begin{equation}\tag{12.17}
\Lambda/\lambda\le\gamma.
\end{equation}
After division by $\lambda$, assume $\lambda=1$, $\Lambda\le\gamma$, and
$|f|\le\mu$.  With $p=u_x$ and $q=u_y$, equation (12.15) becomes
\begin{equation}\tag{12.18}
\frac acp_x+\frac{2b}cp_y+q_y=\frac fc,
\qquad p_y=q_x.
\end{equation}
It follows that
\begin{equation}\tag{12.19}
|Dp|^2+|Dq|^2
\le(1+\gamma)J+\frac12(1+\gamma)\mu(|p_x|+|q_y|),
\qquad
J=q_xp_y-q_yp_x,
\end{equation}
and hence
\begin{equation}\tag{12.20}
|Dp|^2+|Dq|^2
\le2(1+\gamma)J+\frac12(1+\gamma)^2\mu^2.
\end{equation}
Thus $w=p-iq$ is $(K,K')$-quasiconformal with
\begin{equation}\tag{12.21}
K=1+\gamma,
\qquad K'=\frac12(1+\gamma)^2\mu^2.
\end{equation}
When $f=0$, one may take $K=(1+\gamma)/2$; the optimal elementary bound is
$K=(\gamma+\gamma^{-1})/2$.

Write $d_z=\operatorname{dist}(z,\partial\Omega)$ and
$d_{1,2}=\min(d_{z_1},d_{z_2})$, and set
\[
[u]^*_{1,\alpha}
=\sup_{z_1,z_2\in\Omega}
d_{1,2}^{1+\alpha}
\frac{|Du(z_2)-Du(z_1)|}{|z_2-z_1|^\alpha},
\qquad
\|f/\lambda\|_0^{(2)}
=\sup_{z\in\Omega}d_z^2|f/\lambda|.
\]

\begin{theorem}[Interior H\"older estimate for gradients in dimension two]
\label{gt-ch12-thm-interior-holder-gradient-linear}
\dcref{ch12:12.4}\group{chapter-12}\source{gilbarg-trudinger:page-0314}
Let $u\in C^2(\Omega)$ be a bounded solution of (12.15), where
$L$ satisfies (12.16)--(12.17).  For some
$\alpha=\alpha(\gamma)>0$,
\begin{equation}\tag{12.22}
[u]^*_{1,\alpha}
\le C(\gamma)
\left(\|u\|_0+\|f/\lambda\|_0^{(2)}\right).
\end{equation}
\end{theorem}
\begin{proof}
\sketch For $z_1,z_2\in\Omega$, take
$2d=d_{1,2}$ and apply
Theorem~\ref{gt-ch12-thm-holder-quasiconformal-mappings} to
$u_x-iu_y$ on the region at distance greater than $d$ from the boundary.
This bounds the weighted H\"older seminorm by
$C([u]^*_1+\|f/\lambda\|_0^{(2)})$.  The interpolation inequality
\begin{equation}\tag{12.23}
[u]^*_1\le\varepsilon[u]^*_{1,\alpha}+C_1(\varepsilon)|u|_0
\end{equation}
absorbs $[u]^*_{1,\alpha}$ and gives (12.22).
\end{proof}

The corresponding norm estimate is
\begin{equation}\tag{12.24}
|u|^*_{1,\alpha}
\le C(\gamma)\left(|u|_0+|f/\lambda|_0^{(2)}\right).
\end{equation}
If $\Omega$ is $C^2$, $u\in C^2(\overline\Omega)$, and $u=0$ on
$\partial\Omega$, boundary flattening, reflection, and interpolation give
\[
|u|_{1,\alpha;\Omega}
\le C(\gamma,\Omega,|u|_0,|f/\lambda_0|_0),
\]
where $\lambda_0=\inf_\Omega\lambda$.  For
$u=\varphi\in C^2(\overline\Omega)$ on the boundary, the right-hand side
depends instead on $\gamma,\Omega,|\varphi|_2$, and $|f/\lambda_0|_0$.
The same global conclusion holds for
$u\in C^2(\Omega)\cap C^0(\overline\Omega)$ when
$\partial\Omega\in C^{2,\beta}$, $\varphi\in C^{2,\beta}(\overline\Omega)$,
and the coefficients and $f$ lie in $C^\beta(\Omega)$, by smooth
approximation and interior Schauder compactness.

\section*{12.3. The Dirichlet Problem for Uniformly Elliptic Equations}

Consider
\begin{equation}\tag{12.25}
\begin{aligned}
Qu={}&a(x,y,u,u_x,u_y)u_{xx}
+2b(x,y,u,u_x,u_y)u_{xy}\\
&+c(x,y,u,u_x,u_y)u_{yy}
+f(x,y,u,u_x,u_y).
\end{aligned}
\end{equation}
Assume, for some $\beta\in(0,1)$:
\begin{enumerate}
\item[(i)] $a,b,c,f\in
C^\beta(\Omega\times\mathbb R\times\mathbb R^2)$;
\item[(ii)] $Q$ is elliptic and its eigenvalues satisfy
\begin{equation}\tag{12.26}
1\le\Lambda/\lambda\le\gamma(|u|),
\end{equation}
where $\gamma$ is nondecreasing;
\item[(iii)] for a nondecreasing $\mu$ and a constant $\nu\ge0$,
\begin{equation}
|f|/\lambda\le\mu(|u|)(1+|p|+|q|).
\tag{12.27}
\end{equation}
\begin{equation}
(f/\lambda)\operatorname{sign}u
\le\nu(1+|p|+|q|).
\tag{12.28}
\end{equation}
\end{enumerate}

\begin{theorem}[Dirichlet problem for uniformly elliptic equations in dimension two]
\label{gt-ch12-thm-dirichlet-uniformly-elliptic}
\dcref{ch12:12.5}\group{chapter-12}\source{gilbarg-trudinger:page-0317}
Let $\Omega\subset\mathbb R^2$ satisfy an exterior sphere condition and let
$\varphi\in C^0(\partial\Omega)$.  If $Q$ satisfies (i)--(iii), then
\begin{equation}\tag{12.29}
Qu=0\quad\text{in }\Omega,
\qquad u=\varphi\quad\text{on }\partial\Omega
\end{equation}
has a solution
$u\in C^{2,\beta}(\Omega)\cap C^0(\overline\Omega)$.
\end{theorem}
\begin{proof}
\sketch First replace (12.27) by
\begin{equation}\tag{12.30}
|f|/\lambda\le\mu<\infty.
\end{equation}
For $v$ in a bounded subset of a weighted $C^{1,\alpha}$ space, freeze the
coefficients at $(v,Dv)$ and solve
\begin{equation}\tag{12.31}
\bar a u_{xx}+2\bar b u_{xy}+\bar c u_{yy}+\bar f=0
\quad\text{in }\Omega,
\qquad u=\varphi\quad\text{on }\partial\Omega.
\end{equation}
The maximum principle gives
$|u|_0\le\sup_{\partial\Omega}|\varphi|+C\mu=:M_0$, and
Theorem~\ref{gt-ch12-thm-interior-holder-gradient-linear} gives
\begin{equation}\tag{12.32}
|u|^*_{1,\alpha}
\le C\left(|u|_0+\mu(\operatorname{diam}\Omega)^2\right)
\le K,
\qquad
\alpha=\alpha(\gamma(M_0)).
\end{equation}
Define the convex set
\[
\mathcal G=\{v:|v|^*_{1,\alpha}\le K,\ |v|_0\le M_0\}
\]
and let $Tv$ be the solution of (12.31).  Exterior-sphere barriers give,
uniformly in $v$,
\begin{equation}\tag{12.33}
|Tv(z)-\varphi(z_0)|
\le\varepsilon+k_\varepsilon w(z)
\qquad(z_0\in\partial\Omega).
\end{equation}
Thus $T\mathcal G$ is precompact in the weaker weighted $C^1$ norm.
Interior Schauder compactness, (12.33), and uniqueness make $T$ continuous,
so the Schauder fixed-point theorem gives a solution under (12.30).

For the general case divide by $\lambda$ and rewrite (12.27)--(12.28) as
\begin{equation}
|f|\le\mu(|u|)(1+|p|+|q|).
\tag{12.34}
\end{equation}
\begin{equation}
f\operatorname{sign}u\le\nu(1+|p|+|q|).
\tag{12.35}
\end{equation}
Let $\psi_N(t)=t$ for $|t|\le N$ and
$\psi_N(t)=N\operatorname{sign}t$ otherwise, and truncate $f$ in all three
variables $u,p,q$.  Solve
\begin{equation}\tag{12.36}
\begin{aligned}
Q_Nu={}&a(x,y,u,Du)u_{xx}+2b(x,y,u,Du)u_{xy}\\
&+c(x,y,u,Du)u_{yy}+f_N(x,y,u,Du)=0,
\qquad u=\varphi\text{ on }\partial\Omega.
\end{aligned}
\end{equation}
The sign condition gives the $N$-independent bound
\begin{equation}\tag{12.37}
\sup_\Omega|u_N|
\le\sup_{\partial\Omega}|\varphi|
+C_1(\nu,\operatorname{diam}\Omega)=:M.
\end{equation}
The gradient estimate first gives
\begin{equation}\tag{12.38}
[u_N]^*_{1,\alpha}
\le C\left(|u_N|_0+|f_N|_0^{(2)}\right),
\end{equation}
and then, using (12.34) and interpolation,
\begin{equation}\tag{12.39}
|u_N|^*_{1,\alpha}
\le C(M,\gamma(M),\mu(M),\operatorname{diam}\Omega).
\end{equation}
Interior Schauder compactness supplies a subsequential solution of $Qu=0$;
the same uniform barriers as in (12.33) pass the boundary values to the
limit.
\end{proof}

If $\partial\Omega$ and $\varphi$ are in $C^{2,\beta}$ and
$a,b,c,f\in C^\beta(\overline\Omega\times\mathbb R\times\mathbb R^2)$, the
solution is in $C^{2,\beta}(\overline\Omega)$.  Condition (12.28) may be
dropped when a uniform zeroth-order bound is otherwise known, and (12.27)
may be dropped when a uniform gradient bound is known.  The exterior sphere
condition may be replaced by any hypothesis providing the uniform linear
barriers, including an exterior cone condition.

\section*{12.4. Non-Uniformly Elliptic Equations}

Consider the elliptic equation
\begin{equation}\tag{12.40}
Qu=a(x,y,u,u_x,u_y)u_{xx}
+2b(x,y,u,u_x,u_y)u_{xy}
+c(x,y,u,u_x,u_y)u_{yy}=0.
\end{equation}

\begin{definition}[Bounded slope condition]
\label{gt-ch12-def-bounded-slope}
\dcref{ch12:12.41}
\group{chapter-12}\source{gilbarg-trudinger:page-0322}
Let $\Omega\subset\mathbb R^2$ be bounded and
$\varphi:\partial\Omega\to\mathbb R$.  The boundary graph
$\Gamma=(\partial\Omega,\varphi)$ satisfies the bounded slope condition with
constant $K$ if, for every $z_0\in\partial\Omega$, there are affine functions
\[
\pi_{z_0}^\pm(z)
=a^\pm(z_0)\cdot(z-z_0)+\varphi(z_0)
\]
such that
\begin{equation}\tag{12.41}
\pi_{z_0}^-(z)\le\varphi(z)\le\pi_{z_0}^+(z)
\quad(z\in\partial\Omega),
\qquad
|a^\pm(z_0)|\le K.
\end{equation}
\end{definition}

\begin{definition}[Three-point condition]
\label{gt-ch12-def-three-point}
\dcref{ch12:12.41}
\group{chapter-12}\source{gilbarg-trudinger:page-0322}
If $\Omega$ is bounded and convex, $\Gamma$ satisfies the three-point
condition with constant $K$ when every three distinct points of $\Gamma$ lie
in a plane of slope at most $K$.
\end{definition}

The bounded slope condition makes $\varphi$ continuous.  Unless $\Gamma$ is
planar, it also forces $\Omega$ to be convex.  On a straight boundary
segment, $\varphi$ must be affine.  If $\Omega$ is strictly convex, the
three-point condition is equivalent to requiring every plane meeting
$\Gamma$ in at least three points to have slope at most $K$.  In particular,
if $\partial\Omega,\varphi\in C^2$ and the curvature of $\partial\Omega$ is
everywhere positive, then $\varphi$ satisfies a bounded slope condition; its
constant depends only on the minimum curvature and the first two derivatives
of $\varphi$.

\begin{lemma}[Gradient bound from the bounded slope condition]
\label{gt-ch12-lem-bounded-slope-gradient}
\dcref{ch12:12.6}\group{chapter-12}\source{gilbarg-trudinger:page-0323}
Let $\Omega\subset\mathbb R^2$ be bounded, and let
$\varphi$ satisfy a bounded slope condition with constant $K$.  Suppose
$u\in C^2(\Omega)\cap C^0(\overline\Omega)$ satisfies
\begin{equation}\tag{12.42}
Lu=au_{xx}+2bu_{xy}+cu_{yy}=0
\quad\text{in }\Omega,
\qquad u=\varphi\quad\text{on }\partial\Omega,
\end{equation}
where $L$ is elliptic.  Then
\begin{equation}\tag{12.43}
\sup_\Omega|Du|\le K.
\end{equation}
\end{lemma}
\begin{proof}
\sketch If $\varphi$ is affine, uniqueness gives the claim.  Otherwise
$\Omega$ is convex.  Ellipticity and (12.42) imply
$u_{xx}u_{yy}-u_{xy}^2\le0$, with equality only where $D^2u=0$.
At a point with $D^2u\ne0$, subtract the tangent plane.  Its zero set has
four alternating nodal sectors; the maximum principle forces their
continuations to meet $\partial\Omega$ in at least four points.  The
three-point condition therefore bounds the tangent-plane slope by $K$.
Continuity extends this bound to limit points, while $u$ is affine on every
open component where $D^2u=0$.
\end{proof}

\begin{theorem}[Dirichlet problem for non-uniformly elliptic equations]
\label{gt-ch12-thm-dirichlet-nonuniformly-elliptic}
\dcref{ch12:12.7}\group{chapter-12}\source{gilbarg-trudinger:page-0324}
Let $\Omega\subset\mathbb R^2$ be bounded.  Suppose (12.40) is elliptic and
$a,b,c\in C^\beta(\Omega\times\mathbb R\times\mathbb R^2)$ for some
$\beta\in(0,1)$.  If $\varphi$ satisfies a bounded slope condition with
constant $K$, then
\[
Qu=0\quad\text{in }\Omega,
\qquad u=\varphi\quad\text{on }\partial\Omega
\]
has a solution
$u\in C^{2,\beta}(\Omega)\cap C^0(\overline\Omega)$ satisfying
$\sup_\Omega|Du|\le K$.
\end{theorem}
\begin{proof}
\sketch Unless the boundary data are affine, $\Omega$ is convex.  Divide
the coefficients by their largest eigenvalue and solve the uniformly
elliptic regularizations
\begin{equation}\tag{12.44}
Q_\varepsilon u=Qu+\varepsilon\Delta u=0,
\qquad \varepsilon>0.
\end{equation}
Theorem~\ref{gt-ch12-thm-dirichlet-uniformly-elliptic} gives
$u_\varepsilon\in C^{2,\beta}(\Omega)\cap C^0(\overline\Omega)$, and
Lemma~\ref{gt-ch12-lem-bounded-slope-gradient} gives
\begin{equation}\tag{12.45}
\sup_\Omega|Du_\varepsilon|\le K.
\end{equation}
The maximum principle also bounds $|u_\varepsilon|$.  On each
$\Omega'\Subset\Omega$, compactness of
$(x,y,u_\varepsilon,Du_\varepsilon)$ supplies a positive lower ellipticity
bound for the linear equations
\begin{equation}\tag{12.46}
a_\varepsilon u_{xx}
+2b_\varepsilon u_{xy}
+c_\varepsilon u_{yy}
+\varepsilon\Delta u=0.
\end{equation}
Theorem~\ref{gt-ch12-thm-interior-holder-gradient-linear} gives uniform
local H\"older bounds for $Du_\varepsilon$, after which Schauder estimates
give local $C^{2,\alpha\beta}$ compactness.  A diagonal subsequence converges
to a solution of $Qu=0$; (12.45) makes the convergence uniform up to the
boundary.
\end{proof}

The convexity hypothesis cannot be discarded in general.  For example, the
minimal-surface equation
\begin{equation}\tag{12.47}
(1+u_y^2)u_{xx}-2u_xu_yu_{xy}+(1+u_x^2)u_{yy}=0
\end{equation}
on an annulus has a catenoid solution for a sufficiently small height
difference between its boundary circles, but no solution for sufficiently
large height.  Continuous boundary data alone are also insufficient: failure
may occur even for a circular boundary and arbitrarily smooth coefficients.
On the other hand, if $\Omega$ is convex and
\[
\Lambda/\lambda\le\mu(|u|)(1+p^2+q^2),
\]
then (12.40) is solvable for arbitrary $\varphi\in C^2$, without the bounded
slope condition.  Under the additional hypotheses
$a,b,c\in C^\beta(\overline\Omega\times\mathbb R\times\mathbb R^2)$,
$\partial\Omega\in C^{2,\beta}$, and
$\varphi\in C^{2,\beta}(\partial\Omega)$, the solution in
Theorem~\ref{gt-ch12-thm-dirichlet-nonuniformly-elliptic} belongs to
$C^{2,\beta}(\overline\Omega)$.

\begin{proposition}[Equivalence of bounded slope and three-point conditions]
\label{gt-ch12-prop-bounded-slope-three-point-equivalence}
\dcref{ch12:12.48}
\group{chapter-12}\source{gilbarg-trudinger:page-0326}
For a bounded convex domain $\Omega\subset\mathbb R^2$, the bounded slope
condition and the three-point condition are equivalent, with the same
constant $K$.
\end{proposition}
\begin{proof}
\sketch Assume first the bounded slope condition.  If
$z_1,z_2,z_3$ are noncollinear and the plane through
$(z_i,\varphi(z_i))$ has slope vector $a$, then
\begin{equation}\tag{12.48}
a_i^-\cdot(z_j-z_i)
\le a\cdot(z_j-z_i)
\le a_i^+\cdot(z_j-z_i),
\qquad i,j=1,2,3.
\end{equation}
For some vertex $z_i$, either $a$ or $-a$ is a nonnegative combination of
the edge vectors from $z_i$.  Equation (12.48) then gives
$|a|^2\le K|a|$.

Conversely, fix $A\in\Gamma$ and take shrinking triangles through $A$ whose
planes converge to a plane containing a supporting line $L$ of $\Omega$.
For each $Q\in\Gamma\setminus L$, the plane through $Q$ and $L$ has slope at
most $K$ by the three-point condition.  The pointwise supremum and infimum
of these planes on the half-plane containing $\Omega$ are supporting planes
above and below $\Gamma$ at $A$, both with slope at most $K$.  If $A$ lies
on a straight segment, use the line containing that segment for $L$.
\end{proof}
